% Frozen before sampling random allocations or running Aer. Compile with xelatex.
\documentclass[11pt,a4paper]{article}
\usepackage{kotex}
\usepackage{amsmath,amssymb,booktabs,geometry}
\geometry{margin=25mm}
\title{I01 선검증 5 사전 프로토콜: 무작위 표지 수와 블록 간 배분 드리프트}
\author{RELATIONAL ATLAS}
\date{2026년 10월 4일}
\begin{document}
\maketitle

\section{분리할 두 현상}
같은 목표 표지 분포에서 무작위 추출된 표지 수의 차이는
표지 삭제 통계량의 분산을 키울 수 있지만, 목표 분포가 안정적이면
무조건부 평균을 바꾸지 않는다. 반면 두 블록의 \emph{목표 분포
자체}가 다르면 과학 위상을 바꾸지 않아도 표지 삭제 후 평균이
달라질 수 있다. 이 둘을 같은 실험에서 분리한다. 발전연구 주제와
QPU 실험은 보류한다.

\section{사전 고정한 Aer 설계}
선검증 4b의 $n=3$, $L=7$, $\beta=0.45$, 기준 폭
$\beta_r=0.15$, $E=E_0+0.37E_0^2/3$, 검출 바이어스
1.25 rad, $\tau=0,\pi/2$, 비제약/$P=0$ 상태를 사용한다.
세 블록 $A,B,C$의 \emph{과학 공통 위상은 모두 0}이다.
$A,B$의 목표 표지 분포는 균일 $a_t=1/7$,
$C$는 $a_t=[1+0.35\cos(2\pi t/7)]/7$이다.
각 블록$\times$상태$\times$시계$\times$검출 축에서 독립적으로
총 $N=229376$개 표지를 다항분포에서 뽑는다.
따라서 실제 표지 수는 목표값과 다르고, $A,B$도 서로 다르다.
배분 RNG seed 20261023, Aer seed는 20261024에 실행 순번을
더한 값, CPU workers 6을 고정한다. 3블록$\times$2상태
$\times$2시계$\times$3축$\times$7표지 = 252회로다.

\section{사전 검정 가족}
첫 가족은 표지를 \emph{삭제}한 각 축의 성공수만 쓴다.
$A-B$와 $A-C$ 각각에 대해 두 상태$\times$두 시계 위상을
교차한 여덟 검정이다. 축별 두 블록은 각기 $N$개 독립 시행이고,
표지를 매 시행 독립 추출하므로 목표 분포로 적분한 성공확률의
Bernoulli 시행이다. 두 블록의 목표-혼합 확률이 같다는 null에서
축별 합친 성공수를 조건화한 hypergeometric 재배분을
10000회 사용한다. 가족 Bonferroni 문턱은 0.00625다.
정확 모형에서 $A-B$ 네 조건과 $A-C$의 $P=0$ 두 조건은 null,
$A-C$의 비제약 두 조건은 양수다.

둘째 가족은 같은 counts에서 기록된 표지를 사용해 각 블록을
\emph{균일 목표} $1/7$로 표준화한다. 두 블록의 같은 표지별
확률이 같다는 null 아래, 서로 다른 표지별 시행 수를 보존한
축$\times$표지 hypergeometric 재배분을 10000회 한다.
가족은 같은 여덟 대조이고 문턱도 별도로 0.00625다.
과학 위상을 전혀 바꾸지 않았으므로 이 표준화 가족의 정확
null은 여덟 조건 모두에서 참이다. 한 번의 비검출을 일반적
보정 성공의 증명으로 읽지 않는다.

정확 확률의 사전 pilot은 $A-C$ 비제약의 최대 합산확률 차이를
$\tau=0,\pi/2$에서 약 0.0275/0.0134로 예상했다.
또 균일 목표의 독립 무작위 표지는 고정 표지 수 대비 축별
표지 삭제 분산을 비제약 상태에서 약 1.35/1.16배 키우지만,
$P=0$에서는 1배다. Aer 실행 뒤 결과에 맞춰 샷 예산,
seed 또는 검정 가족을 바꾸지 않는다. 별도 시계 레지스터와
미지 목표 분포 자체의 추정은 이번 실험에 포함하지 않는다.
\end{document}
